% Springer LNCS draft for the CPV301 driver-drowsiness project.
% Compile with the official llncs.cls and splncs04.bst supplied by Springer
% or by the lecturer. Do not replace them with a locally emulated class.
\documentclass[runningheads]{llncs}

\usepackage[T1]{fontenc}
\usepackage[utf8]{inputenc}
\usepackage{graphicx}
\usepackage{amsmath,amssymb}
\usepackage{booktabs}
\usepackage{multirow}
\usepackage{array}
\usepackage[hidelinks]{hyperref}

\graphicspath{{figures/}}
\emergencystretch=1em

% Keeps the draft compilable while figures are being assembled. Every missing
% figure box must be replaced by the real project image before submission.
\newcommand{\projectfigure}[2][]{%
  \IfFileExists{figures/#2}{%
    \includegraphics[#1]{#2}%
  }{%
    \fbox{\parbox[c][35mm][c]{0.90\linewidth}{\centering
      Missing project figure: \texttt{#2}}}%
  }%
}

\title{An Explainable Real-Time Driver Drowsiness Detection System Using Facial Geometry and Temporal Rules}

\titlerunning{Explainable Driver Drowsiness Detection}

\author{Nguyen Bao Dang \and
        Pham Vu Gia Kiet \and
        Nguyen Xuan Binh\thanks{Corresponding author.} \and
        Ha Thuc Khuong Ninh}

\authorrunning{N. B. Dang et al.}

% Confirm the official department and campus wording before submission.
\institute{Department of Artificial Intelligence, FPT University,\\
Ho Chi Minh City, Vietnam\\
\email{ngxbinh9999@gmail.com}}

\begin{document}
\maketitle

\begin{abstract}
Camera-based driver monitoring requires temporal reasoning because normal blinks,
speech-related mouth opening, and brief head movements may resemble fatigue cues
in individual frames. This paper presents an explainable academic prototype that
combines pretrained MediaPipe Face Mesh landmarks with Eye Aspect Ratio (EAR),
Mouth Aspect Ratio (MAR), eye-line roll, and duration rules expressed in seconds.
The method was evaluated on 41,793 DDD images, 1,448 full-face yawn/no-yawn
images, and six self-recorded videos from three participants. On usable static
images, the DDD experiment achieved 0.8701 precision and 0.4673 recall, whereas
yawn detection at MAR $>0.60$ achieved 1.0000 precision and 0.6551 recall. In
video, event recall was 0.8000 for eye closure, 0.6000 for yawning, and 0.5000
for lateral head tilt. Participant-specific results varied substantially because
recording and annotation protocols were not standardized. The results show that
second-based temporal tracking supports videos with different frame rates, but
fixed thresholds, roll-only head geometry, and limited participant coverage
prevent safety-critical deployment.

\keywords{Driver drowsiness detection \and Facial landmarks \and Eye Aspect Ratio
\and MediaPipe Face Mesh \and Temporal reasoning}
\end{abstract}

\section{Introduction}

Driver drowsiness is a safety-relevant condition that cannot be inferred reliably
from a single visual observation. A normal blink, a speech-related mouth opening,
or a short head movement can resemble a fatigue cue in an isolated frame. Camera-
based monitoring is nevertheless attractive for an academic prototype because it
is contactless and can produce interpretable evidence from the eyes, mouth, and
head \cite{sahayadhas2012review}.

This work investigates a lightweight alternative to training a project-specific
drowsiness classifier. MediaPipe Face Mesh provides pretrained facial landmarks,
while project-defined geometric measurements and temporal rules determine the
warning state. Measuring persistence in seconds rather than by a fixed number of
frames is important because the submitted videos have different frame rates. The
system is evaluated on public static images and continuous videos, with results
reported at frame, event, aggregate, and participant levels.

The practical contributions are: (i) integration of EAR, MAR, and lateral head
roll in one real-time pipeline; (ii) duration tracking shared by webcam inference
and offline evaluation; (iii) reproducible sample subsets for submission smoke
tests; and (iv) explicit analysis of threshold trade-offs, participant variation,
and failure cases. The system is an explainable course prototype rather than a
certified vehicle-safety device.

\section{Related Work}

Driver-drowsiness monitoring methods are commonly grouped into physiological,
vehicle-based, and behavioural approaches \cite{sahayadhas2012review}. Camera-based behavioural systems can measure several facial cues without wearable
sensors. Sigari et al. combined eyelid, PERCLOS, and head-orientation information
within a driver face-monitoring system \cite{sigari2013monitoring}. Their work
supports multi-cue monitoring but also indicates that a single observable signal
is not sufficient in all conditions.

Soukupov\'a and \v{C}ech introduced the six-landmark EAR formulation for blink
analysis \cite{soukupova2016ear}. Alioua et al. used mouth motion across
consecutive frames for yawning-based fatigue detection \cite{alioua2014yawning}.
Both approaches motivate the combination of interpretable facial geometry with
temporal continuity rather than isolated frame decisions.

MediaPipe provides a low-latency framework for perception pipelines
\cite{lugaresi2019mediapipe}. Its monocular Face Mesh method estimates a dense
facial surface with 468 landmarks \cite{kartynnik2019facemesh}, while Attention
Mesh improves landmark quality in semantically important eye and lip regions
\cite{grishchenko2020attention}. In the present project, these pretrained methods
are used only for landmark extraction; the drowsiness decisions are produced by
explicit ratios, an angle, and duration rules.

Recent systems increasingly combine eye, mouth, and head information. Zhu et al.
used facial video sequences with EAR, MAR, PERCLOS, and landmarks
\cite{zhu2022fatigue}. Albadawi et al. supplied eye, mouth, and head-pose features
to machine-learning classifiers \cite{albadawi2023visual}. Kim et al. combined
landmark-based eye-closure detection with head-pose recognition and temporal
continuity \cite{kim2023monitoring}. The present work does not propose a new
landmark model or claim state-of-the-art performance. Its focus is transparent
integration, second-based timing, and multi-level evaluation of a small academic
prototype.

\section{Proposed Method}

\subsection{System Overview}

\begin{figure}[t]
\centering
\projectfigure[width=0.98\linewidth]{system_pipeline.png}
\caption{Pipeline from image, video, or webcam input through Face Mesh,
geometric cues, temporal rules, real-time warning, and offline evaluation.}
\label{fig:pipeline}
\end{figure}

Figure~\ref{fig:pipeline} summarizes the processing pipeline. Webcam frames,
recorded videos, or static images are converted from BGR to RGB and processed by
MediaPipe Face Mesh. If a face is found, the system computes EAR, MAR, and an
eye-line roll angle. Static images use instantaneous geometric decisions. Webcam
and video modes additionally require each condition to remain continuously true
for its configured duration. A no-face frame is counted during evaluation and
resets active temporal trackers.

\subsection{Geometric Features}

Six landmarks are used for each eye: left
$\{33,160,158,133,153,144\}$ and right
$\{362,385,387,263,373,380\}$. For eye points $p_1,\ldots,p_6$, EAR is

\begin{equation}
\mathrm{EAR}=\frac{\lVert p_2-p_6\rVert_2+\lVert p_3-p_5\rVert_2}
{2\lVert p_1-p_4\rVert_2}.
\label{eq:ear}
\end{equation}

The two eye values are averaged; a smaller EAR indicates greater closure. Mouth
landmarks $\{78,308,13,14\}$ represent the corners and upper/lower lip centres:

\begin{equation}
\mathrm{MAR}=\frac{\lVert p_{\mathrm{upper}}-p_{\mathrm{lower}}\rVert_2}
{\lVert p_{\mathrm{left}}-p_{\mathrm{right}}\rVert_2}.
\label{eq:mar}
\end{equation}

A larger MAR indicates greater mouth opening. Outer eye corners 33 and 263 form
the line used to estimate absolute roll:

\begin{equation}
\theta_{\mathrm{roll}}=
\left|\frac{180}{\pi}\operatorname{atan2}(y_r-y_l,x_r-x_l)\right|.
\label{eq:roll}
\end{equation}

This last cue represents lateral tilt. It does not estimate complete 3D head
pose and can therefore miss forward nodding dominated by pitch.

\subsection{Temporal Decision Rules}

For recorded video, elapsed time increases by $1/\mathrm{FPS}$ for each frame;
for a webcam, monotonic wall-clock time is used. An alert activates only after a
condition remains continuously true for the required duration. The final shared
configuration is shown in Table~\ref{tab:thresholds}.

\begin{table}[t]
\caption{Final geometric thresholds and minimum durations.}
\label{tab:thresholds}
\centering
\small
\begin{tabular}{@{}lccc@{}}
\toprule
Cue & Condition & Duration & Intended behaviour \\
\midrule
Eye closure & EAR $<0.21$ & 1.000 s & Prolonged closure \\
Yawn & MAR $>0.60$ & 0.500 s & Sustained mouth opening \\
Head tilt & $|\theta_{\mathrm{roll}}|>15^{\circ}$ & 0.667 s & Lateral head lean \\
\bottomrule
\end{tabular}
\end{table}

MediaPipe Face Mesh is pretrained. The project does not train or fine-tune a new
classifier; therefore optimizer, loss, learning rate, epoch count, and training
augmentation do not apply to the final rule-based decision stage. The relevant
configurable parameters are the geometric thresholds and temporal durations.

\section{Experimental Setup}

\subsection{Datasets and Submission Samples}

The full Driver Drowsiness Dataset (DDD) contains 41,793 images: 22,348 drowsy
and 19,445 non-drowsy images \cite{nasri_ddd}. Seventeen images had no detected
face, leaving 41,776 usable samples. The full-face part of the yawn/no-yawn
dataset contains 1,448 images \cite{raju_yawn}; two had no detected face,
leaving 1,446 usable samples. Cropped eye-only images were excluded because Face
Mesh requires a full face.

To keep the submitted project reproducible and below the upload limit, the
repository also contains a deterministic 300-image smoke-test subset: 50 images
for each DDD class and 50 images for each yawn/no-yawn class in both the train
and test splits. These samples verify installation and data traversal; the full-
dataset results reported below were not computed from this reduced subset.

\begin{figure}[t]
\centering
\projectfigure[width=0.96\linewidth]{sample_dataset_overview.png}
\caption{Representative images from the submitted sample subsets. The figure
illustrates DDD drowsy/non-drowsy and yawn/no-yawn variation; it is not an
additional quantitative experiment.}
\label{fig:samples}
\end{figure}

\subsection{Recorded Videos}

Three participants each supplied an alert video and an acted-drowsiness video.
The six videos contain 11,938 frames and 21 labeled events: 10 eye closures,
five yawns, and six head tilts. Recording protocols were not standardized;
frame rate, resolution, camera placement, performed behaviour, and annotation
procedure differ across participants. Consequently, participant-specific
results are treated as primary evidence and pooled results as supplementary
summaries. The experiment has $N=3$ and cannot establish population-level
generalization.

\begin{table}[t]
\caption{Submitted-video metadata and labeled events.}
\label{tab:videos}
\centering
\scriptsize
\setlength{\tabcolsep}{3.2pt}
\begin{tabular}{@{}llrrrrrr@{}}
\toprule
Participant & Video & Frames & FPS & No-face & Eye & Yawn & Tilt \\
\midrule
\multirow{2}{*}{Binh} & Alert & 1,351 & 16.6 & 0 & 0 & 0 & 0 \\
 & Drowsy & 1,421 & 16.6 & 21 & 5 & 3 & 3 \\
\multirow{2}{*}{Kiet} & Alert & 2,287 & 29.9 & 0 & 0 & 0 & 0 \\
 & Drowsy & 3,225 & 29.3 & 0 & 4 & 2 & 3 \\
\multirow{2}{*}{Ninh} & Alert & 1,834 & 29.9 & 0 & 0 & 0 & 0 \\
 & Drowsy & 1,820 & 29.9 & 0 & 1 & 0 & 0 \\
\midrule
Total & Six videos & 11,938 & -- & 21 & 10 & 5 & 6 \\
\bottomrule
\end{tabular}
\end{table}

Ground truth is represented by event type, start time, and end time. Frame-level
evaluation compares every prediction with the active labeled interval. An event
alert is matched when it overlaps at least 30\% of a labeled interval. Each
labeled interval can generate one hit; an unmatched label is a miss, and an
unmatched alert is a false alarm.

\subsection{Metrics and Implementation}

Static and frame-level evaluation use accuracy, precision, recall, and F1-score.
Event precision is $H/(H+F)$ and event recall is $H/(H+M)$, where $H$, $F$, and
$M$ denote hits, false alarms, and misses. A metric is reported as not applicable
when the required positive support is absent. Accuracy is interpreted together
with precision, recall, F1, and event counts because negative video frames are
the majority.

The verified interpreter was Python 3.11.9. Package versions were OpenCV
4.11.0, MediaPipe 0.10.21, NumPy 1.26.4, and Matplotlib 3.11.0. Shared
configuration and temporal-logic modules are used by both real-time inference
and evaluation programs.

\section{Results and Discussion}

\subsection{Static-Image Evaluation}

Table~\ref{tab:static} reports the full-dataset results. DDD precision is high,
but recall is 0.4673, so more than half of labeled drowsy images are missed. At
MAR $>0.60$, no false positive occurs in yawn/no-yawn images, while 249 yawns
are missed. Figure~\ref{fig:staticcm} shows that the main weakness in both tasks
is false-negative detection rather than false positives.

\begin{table}[t]
\caption{Full public-image dataset results using the final configuration.}
\label{tab:static}
\centering
\scriptsize
\setlength{\tabcolsep}{3.4pt}
\begin{tabular}{@{}lrrrrrrrr@{}}
\toprule
Dataset & TP & TN & FP & FN & Acc. & Prec. & Recall & F1 \\
\midrule
DDD & 10,438 & 17,878 & 1,559 & 11,901 & .6778 & .8701 & .4673 & .6080 \\
Yawn/no-yawn & 473 & 724 & 0 & 249 & .8278 & 1.0000 & .6551 & .7916 \\
\bottomrule
\end{tabular}
\end{table}

\begin{figure}[t]
\centering
\projectfigure[width=0.96\linewidth]{fig01_static_confusion_matrices.png}
\caption{Row-normalized confusion matrices on the full processed DDD and
yawn/no-yawn datasets. Images without a detected face are excluded from the
matrices and reported separately.}
\label{fig:staticcm}
\end{figure}

Lowering MAR from 0.60 to 0.45 increased yawn recall from 0.6551 to 0.8615 and
F1 from 0.7916 to 0.9256 without a false positive on this static dataset.
Nevertheless, 0.60 was retained as a conservative setting because static
no-yawn images do not adequately represent speech, laughter, or other prolonged
mouth movements that can occur in video.

\subsection{EAR Configuration Analysis}

Four EAR/duration settings were tested on one participant's development videos
(Table~\ref{tab:ear}). All settings detected the five labeled eye-closure events.
Configuration B reduced false alarms from 10 to six by increasing persistence
from 0.67 to 1.00 s. Its frame F1 (0.6126) is slightly below the baseline
(0.6274), so B is selected as a false-alarm/duration trade-off rather than an
absolute optimum. Figure~\ref{fig:ear} visualizes both frame metrics and event
errors. Because the comparison uses one participant, it is a sensitivity study,
not evidence of general superiority.

\begin{table}[t]
\caption{EAR and duration configurations on the development participant.}
\label{tab:ear}
\centering
\small
\begin{tabular}{@{}lccrrrr@{}}
\toprule
Config. & EAR & Duration & Hits & False alarms & Recall & F1 \\
\midrule
Baseline & .21 & .67 s & 5/5 & 10 & .6578 & .6274 \\
A & .18 & .67 s & 5/5 & 7 & .4702 & .5338 \\
B (selected) & .21 & 1.00 s & 5/5 & 6 & .5916 & .6126 \\
C & .19 & .83 s & 5/5 & 6 & .4834 & .5455 \\
\bottomrule
\end{tabular}
\end{table}

\begin{figure}[t]
\centering
\projectfigure[width=0.96\linewidth]{fig04_ear_configuration_comparison.png}
\caption{Eye-closure frame metrics and event errors for four configurations.
B* denotes the selected false-alarm/duration trade-off.}
\label{fig:ear}
\end{figure}

\subsection{Video Results}

Across all videos, the system detected eight of 10 eye closures, three of five
yawns, and three of six head tilts (Table~\ref{tab:event}). Eye closure has the
highest event recall but the lowest event precision because ten unmatched eye
alerts were produced. Yawn detection is more conservative, while head-tilt
precision and recall are both 0.50.

\begin{table}[t]
\caption{Supplementary pooled event-level video results.}
\label{tab:event}
\centering
\small
\begin{tabular}{@{}lrrrrrr@{}}
\toprule
Event & GT & Hit & Miss & False alarm & Precision & Recall \\
\midrule
Eye closure & 10 & 8 & 2 & 10 & .4444 & .8000 \\
Yawn & 5 & 3 & 2 & 1 & .7500 & .6000 \\
Head tilt & 6 & 3 & 3 & 3 & .5000 & .5000 \\
\bottomrule
\end{tabular}
\end{table}

Participant-specific event metrics in Fig.~\ref{fig:subject} provide the more
informative comparison. Binh achieved complete eye/yawn recall but no head-tilt
hit; Kiet achieved complete head-tilt recall but detected only half of eye
closures and none of the labeled yawns; Ninh contributed one detected eye event
and had no positive yawn or tilt support. An N/A cell is unsupported or undefined
and must not be interpreted as a zero score.

\begin{figure}[t]
\centering
\projectfigure[width=0.92\linewidth]{fig06_video_event_metrics_per_subject.png}
\caption{Primary event-level recall and precision by participant. Differences
reflect both system behaviour and non-standardized recording/annotation
protocols; N/A denotes an unsupported or undefined metric.}
\label{fig:subject}
\end{figure}

Pooled frame-level F1 was 0.8398 for eye closure, 0.5118 for yawning, and 0.5507
for head tilt. The corresponding recall values were 0.7936, 0.3453, and 0.3964.
Accuracy exceeded 0.92 for every cue because negative frames dominate, so it
does not by itself establish reliable warning performance.

\subsection{Failure Cases}

The quantitative results identify several failure modes. DDD recall of 0.4673
shows that a fixed EAR threshold misses many labeled drowsy images. Eye closure
produced ten event-level false alarms, which may be associated with squinting,
landmark jitter, annotation timing, or a common threshold applied to different
faces. MAR $=0.60$ reduces false positives but misses mouth openings that are
not large or persistent enough. Roll-only geometry misses forward nodding when
pitch changes without a large lateral eye-line angle. Finally, the Binh drowsy
video contains 21 no-face frames (1.48\%); each loss resets temporal state and
can fragment a genuine event.

These explanations are hypotheses consistent with the measurements rather than
controlled causal findings, because several recording factors changed at the
same time. Before final submission, this subsection should include three to five
authentic frames with video filename and timestamp: at least one eye false
negative, one eye false alarm, one missed yawn, one forward-pitch miss, and one
no-face example. No failure-case image should be staged or fabricated.

\section{Limitations and Future Work}

The video experiment contains only three participants and uses acted behaviours.
FPS, resolution, camera placement, event performance, and annotation procedure
were not standardized. Fixed thresholds may not transfer across facial geometry,
eyewear, illumination, and camera viewpoints. Static image labels also do not
represent the temporal context of a real warning. The method estimates roll but
not full pitch/yaw, and loss of facial landmarks interrupts duration tracking.
No learned baseline was evaluated under the same protocol.

Future work should standardize capture and annotation, recruit more participants,
separate development and test participants, calibrate neutral EAR/MAR baselines
per user, estimate full 3D head pose, and include natural speech and laughter as
negative yawn examples. With a larger participant sample, evaluation should
report participant-level macro averages and uncertainty rather than relying on
pooled frames.

\section{Conclusion}

This paper presented an explainable driver-drowsiness prototype using pretrained
facial landmarks, EAR, MAR, head roll, and duration measured in seconds. The
system supports webcam inference and offline evaluation without training a new
drowsiness classifier. Static results showed high precision but limited recall,
and video evaluation showed that eye closure was detected more reliably than
yawning or lateral head tilt. Increasing eye persistence reduced false alarms on
the development participant, but performance varied strongly across the three
participants. The current evidence supports a reproducible academic prototype,
not deployment as a safety-critical driver-monitoring system.

% Add only verified statements before final submission, for example:
% \section*{Acknowledgements}
% \section*{Author Contributions}
% \section*{Funding and Competing Interests}
% Do not claim funding status, conflicts, ethics approval, or contribution roles
% until every author has confirmed the wording.

\bibliographystyle{splncs04}
\bibliography{references}

\end{document}
